\section{Introduction}

Let $G60$ denote the connected quartic graph on sixty vertices identified
with $\mathrm{AT4val}[60,6]$. The purpose of this paper is to determine
its full automorphism group, describe that group abstractly, and provide
an explicit certificate connecting the native permutation action to the
abstract model.

The graph admits a canonical fixed-point-free involution $a$. Its orbits
form thirty two-element fibers, and the corresponding quotient is a
thirty-vertex graph denoted $G30$. This quotient provides the natural
starting point for the automorphism calculation. Every automorphism of
$G30$ has exactly two lifts to $G60$, producing a subgroup of order
$480$ in $\operatorname{Aut}(G60)$.

The existence of a subgroup of order $480$ does not by itself determine
the full automorphism group. A priori, the native graph could possess
additional symmetries that do not preserve the chosen fiber system. We
therefore perform an independent exhaustive enumeration directly on the
sixty-vertex graph. The enumeration returns exactly $480$
automorphisms. Moreover, every enumerated automorphism preserves the
$a$-fiber partition and centralizes the deck involution. It follows that
the lifted subgroup is the full automorphism group.

The order alone does not identify the group. We therefore calculate its
center, derived series, abelianization, element-order profile, and
index-two subgroups. These invariants lead to the fiber-product model
\[
S_5 \times_{C_2} D_8
=
\left\{
(\sigma,d)\in S_5\times D_8:
\operatorname{sgn}(\sigma)=\chi(d)
\right\},
\]
where $\chi\colon D_8\to C_2$ is the quotient character whose kernel is
a Klein four subgroup. We then construct a generator correspondence and
extend it to an explicit bijection between all $480$ elements of the
native automorphism group and all $480$ elements of the fiber-product
model. The bijection is verified against every ordered product.

The principal result is therefore the following.

\begin{theorem}
For the native graph $G60\cong \mathrm{AT4val}[60,6]$,
\[
\left|\operatorname{Aut}(G60)\right|=480
\]
and
\[
\operatorname{Aut}(G60)
\cong
S_5\times_{C_2}D_8,
\]
where the fiber product matches the sign character of $S_5$ with the
quotient character of $D_8$ having Klein four kernel.
\end{theorem}

This manuscript concerns the sixty-vertex graph $G60$. It is distinct
from the companion study of the thirty-vertex graph
\[
X=\operatorname{KC}(L(P)),
\]
whose automorphism group has order $720$ and is isomorphic to
$S_5\times S_3$. The two graphs and their automorphism groups must not be
conflated.

The paper is organized as follows. Section~2 fixes the native graph and
its exact data model. Section~3 develops the canonical quotient and the
lifting construction. Section~4 proves completeness by independent
native enumeration. Section~5 determines the internal group structure.
Section~6 identifies the fiber-product model. Section~7 constructs the
explicit isomorphism. Section~8 records the computational certificates,
and Section~9 summarizes the classification.
